\item Obviously, buyers would prefer general warranty deeds, while sellers would
prefer to make no promises. Why would a seller agree to a warranty deed? Why
would a buyer agree to a special warranty deed? And when would a quitclaim deed
be useful? Put another way, can you imagine a situation where a seller
could not guarantee ownership of the land being sold?

\item Identify the covenant(s) of title, if any, violated based on the following
fact patterns:
\begin{itemize}
\item Alice sells to Bob with a general warranty deed. Alice doesn't own the
house.

\item Charlie and Dave are married and own a house together. Without telling
Dave, Charlie sells his share of the house to Eve with a special warranty deed.
\having{us-v-craft-ii}{Charlie and Dave own the house as tenants by the
entirety.}{Ordinarily there's nothing wrong with selling one's share of
ownership, but Charlie and Dave own the house in a manner which prevents Charlie
from selling his share unilaterally.\footnote{Specifically, they
own the house as tenants by the entirety, discussed later in
\mref{us-v-craft-ii}.}}{Ordinarily there's nothing wrong with selling one's
share of ownership, but Charlie and Dave own the house in a manner which
requires Dave's consent.\footnote{Specifically, they own the house using a
property device called  ``tenancy by the entirety.''}}

\item Frank owns a house. \having{intro-easements}{Frank gives Grace an
easement to walk across the front lawn.}{Frank gives Grace an irrevocable
property right to walk across the front lawn.\footnote{Such a right, called an
``easement,'' is discussed in \mref{intro-easements}.}}{Frank gives Grace a
type of irrevocable property right, called an ``easement,'' to walk across the
front lawn.} After Frank dies, Heidi inherits the house. Not knowing of Grace's
rights, Heidi sells the house to Ivan using a general warranty deed. (What about
a special warranty deed?)
\end{itemize}
